% ch1_basics.tex — 复数基础与解析函数

\subsection{复数域 $\mathbb{C}$}

复数域 $\mathbb{C}$ 是在 $\mathbb{R}$ 上通过添加虚数单位 $i$（满足 $i^2 = -1$）所得的扩域。
从代数角度看，$\mathbb{C} = \mathbb{R}[x]/(x^2+1)$。复数域是\textbf{代数封闭的}
\eng{algebraically closed}：每个非常数复系数多项式在 $\mathbb{C}$ 中都有根。

\begin{remark}
$\mathbb{C}$ 不能赋予全序使其成为有序域 \eng{ordered field}（若 $i > 0$，则 $i^2 = -1 > 0$，矛盾）。
\end{remark}

\subsection{复数的表示}

每个复数 $z \in \mathbb{C}$ 可唯一地写成 $z = x + iy$（$x,y \in \mathbb{R}$），或用极坐标
\[
  z = r e^{i\theta}, \quad r = |z|,\; \theta = \arg z.
\]

\begin{definition}[辐角 \eng{argument}]
$\arg z$ 是多值函数，主值分支 \eng{principal branch} 定义在割平面
$\mathbb{C}^- = \mathbb{C} \setminus (-\infty, 0]$ 上取 $(-\pi, \pi)$。
在 $\mathbb{C}^-$ 上，$\arg$ 是连续的。
\end{definition}

\subsection{指数函数与三角函数}

通过幂级数定义：
\[
  e^z = \sum_{n=0}^{\infty} \frac{z^n}{n!}, \quad
  \cos z = \frac{e^{iz} + e^{-iz}}{2}, \quad
  \sin z = \frac{e^{iz} - e^{-iz}}{2i}.
\]

\begin{lemma}[指数函数的基本性质]
$e^z$ 满足：
\begin{enumerate}  \item $e^{z_1 + z_2} = e^{z_1} e^{z_2}$（加法公式）；
  \item $e^z$ 以 $2\pi i$ 为周期；
  \item $e^z : \mathbb{C} \to \mathbb{C} \setminus \{0\}$ 是满射。
\end{enumerate}
\end{lemma}

\begin{remark}
三角函数与双曲函数 \eng{hyperbolic functions} 通过 $\cosh z = \cos(iz)$, $\sinh z = -i\sin(iz)$ 联系。
映射 $z \mapsto \cos z$ 将竖直线映为双曲线，水平线映为椭圆。
\end{remark}

\subsection{解析函数}

\begin{definition}[解析 \eng{analytic} / 全纯 \eng{holomorphic}]
设 $U \subset \mathbb{C}$ 为开集，$f: U \to \mathbb{C}$。若 $f$ 在 $U$ 上处处复可微
\eng{complex differentiable}，即极限
\[
  f'(z_0) = \lim_{h \to 0} \frac{f(z_0+h) - f(z_0)}{h}
\]
在每点 $z_0 \in U$ 存在，则称 $f$ 在 $U$ 上解析（全纯）。

若 $U = \mathbb{C}$，则称 $f$ 为\textbf{整函数} \eng{entire function}。
\end{definition}

\begin{remark}
复可微比实可微条件强得多：复可微蕴含无穷次可微，甚至可展开为幂级数。
这是复分析区别于实分析的核心。
\end{remark}
